\chapter{The SSH Terminal Console}
\label{ch:ssh}
\section{The architecture}
The console is not a feature of \cmd{kbd}. It is a \emph{second, independent} \cmd{sshd} that exists only to
authenticate an operator and force one program. Host keys, \file{authorized_keys} and PTY handling all live
in that \cmd{sshd}.

\begin{figure}[ht]
\centering
\begin{tikzpicture}[x=1cm,y=1cm,scale=0.9,transform shape]
  \node[box,minimum width=2.4cm] (cl) at (0,0) {operator\\[-1pt]\scriptsize \texttt{ssh -t -p 2222}};
  \node[box,minimum width=3.2cm] (sd) at (4.4,0) {\texttt{sshd@kb-operator}\\[-1pt]\scriptsize key auth, \texttt{Match User operator}};
  \node[box,minimum width=3.6cm] (wr) at (9.2,0) {\texttt{kb-tui-session-wrapper.sh}\\[-1pt]\scriptsize ForceCommand};
  \node[box,minimum width=3.2cm] (s1) at (9.2,-2.1) {\texttt{kbctl ssh session-start}\\[-1pt]\scriptsize audit callback, best-effort};
  \node[dark,minimum width=3.2cm] (tui) at (9.2,-4.2) {\texttt{kbopd} (TUI on the PTY)};
  \node[box,minimum width=3.2cm] (s2) at (9.2,-6.3) {\texttt{kbctl ssh session-end}\\[-1pt]\scriptsize audit callback};
  \node[box,minimum width=2.8cm] (kbd) at (3.4,-4.2) {\texttt{kbd}\\[-1pt]\scriptsize \texttt{kba.sock} gRPC};
  \draw[flow] (cl) -- (sd); \draw[flow] (sd) -- (wr);
  \draw[flow] (wr) -- node[lbl,right]{1} (s1); \draw[flow] (s1) -- node[lbl,right]{2} (tui);
  \draw[flow] (tui) -- node[lbl,right]{3, on quit} (s2);
  \draw[flow] (tui) -- node[lbl,above]{gRPC} (kbd);
  \draw[flowd] (s1.west) -| node[lbl,pos=0.25,above]{record} (kbd.north);
\end{tikzpicture}
\caption{One SSH session. The wrapper records the start, runs the TUI, records the end, and returns the
TUI's exit status to the client.}
\end{figure}
Both audit callbacks are \textbf{best-effort by design}. If \cmd{kbd} is down, \cmd{kbctl ssh session-start}
prints a warning to the operator's terminal and exits 0; an audit-logging outage must never become a denial of
access to the very tool you need during an incident.

\section{Commands}
\begin{lstlisting}[style=console]
$ sudo kbctl ssh up            # enable + start sshd@kb-operator, wait for :2222
$ kbctl ssh status             # exit 0 only if active AND port answering
$ sudo kbctl ssh add-key ~/.ssh/id_ed25519.pub
key authorized for operator
$ kbctl ssh connect            # to 127.0.0.1:2222 as operator, starting it if needed
$ kbctl ssh connect kb-host    # to a remote host
$ kbctl ssh connect kb-host -i ~/.ssh/ops_key --port 2222
$ sudo kbctl ssh down          # stop + disable
\end{lstlisting}
\cmd{ssh connect} replaces itself with \cmd{ssh -t} (so your terminal, window resizes and exit status are
real SSH's). For loopback targets it adds \texttt{UserKnownHostsFile=/dev/null} and
\texttt{StrictHostKeyChecking=no}, because the host key is regenerated per machine and would otherwise
pollute \file{known_hosts}; \emph{remote} hosts keep normal host-key checking.

\subsection{\cmd{add-key} in detail}
It accepts a public-key file or \texttt{-} for stdin, requires \textbf{exactly one line}, checks the key type
prefix (\texttt{ssh-ed25519}, \texttt{ssh-rsa}, \texttt{ecdsa-sha2-*}, \texttt{sk-*}), skips a key already present,
and appends to \file{/etc/kb/authorized_keys} with \cmd{tee -a} under \cmd{sudo} (the file is \texttt{operator:kb} 0600).
\begin{lstlisting}[style=console]
$ ssh-add -L | head -1 | sudo kbctl ssh add-key -     # from an agent
$ printf 'garbage\n' | kbctl ssh add-key -
error: not an OpenSSH public key (want ssh-ed25519/ecdsa-sha2-*/sk-*/ssh-rsa ...)
\end{lstlisting}

\section{The sshd configuration, line by line}
\lstinputlisting[style=conf,caption={\file{/etc/kb/sshd_kb-operator.conf} (packaged)}]{../../packaging/ssh/kb-operator.conf}
\begin{center}\small
\begin{tabularx}{\textwidth}{@{}L{4.6cm}Y@{}}
\toprule
\textbf{Directive} & \textbf{Effect}\\
\midrule
\texttt{Port 2222} & Dedicated port, separate from the host's port 22.\\
\texttt{HostKey /etc/kb/ssh/...} & Host key private to this instance.\\
\texttt{Match User operator} & Everything below applies only to the \texttt{operator} account.\\
\texttt{ForceCommand ...wrapper.sh} & Whatever the client asks to run, this runs instead.\\
\texttt{AuthorizedKeysFile /etc/kb/authorized\_keys} & Keys live outside any home directory.\\
\texttt{X11Forwarding no}, \texttt{AllowTcpForwarding no} & No tunnels through the console.\\
\texttt{PermitTTY yes} & The TUI needs a PTY.\\
\bottomrule
\end{tabularx}
\end{center}
\begin{kbnote}[HARDENING YOU MAY WANT TO ADD]
The packaged file does not disable password authentication for non-\texttt{operator} users. The
\texttt{operator} account has no password, so it is moot for that account, but adding
\texttt{PasswordAuthentication no} and \texttt{KbdInteractiveAuthentication no} makes the whole instance
key-only. Edit \file{/etc/kb/sshd_kb-operator.conf}; it is a protected config file and survives upgrades.
\end{kbnote}

\section{Check the console on your host}
\label{sec:nologin}
OpenSSH runs a \texttt{ForceCommand} through the account's login shell, which a \texttt{nologin} shell refuses.
\file{provision.sh} therefore creates \texttt{operator} with \file{/bin/sh}, and the package's migration step fixes an
existing account that still has \texttt{nologin} or \texttt{false}. This has \textbf{not} been exercised on a live host (no root in
the build sessions), so run this check once after \cmd{kbctl up}:
\resetsteps
\step{Connect}
\begin{lstlisting}[style=console]
$ kbctl ssh connect
\end{lstlisting}
\step{Read the result}
\begin{itemize}
\item \textbf{The TUI appears.} Nothing to do.
\item \textbf{``This account is currently not available.''} and the connection closes: the account still has a \texttt{nologin} shell (for example the package was installed before this fix and the migration did not run). Fix it as below.
\end{itemize}
\step{If refused, give the account a plain shell by hand}
\begin{lstlisting}[style=console]
$ sudo usermod --shell /bin/sh operator
$ kbctl ssh connect
\end{lstlisting}
This is safe in this design: the only way in is a key in \file{/etc/kb/authorized_keys}, the config forces
the wrapper and forbids forwarding, and the account has no password.

\section{Using the TUI}
\subsection{The four tabs}
\texttt{Tab} moves forward, \texttt{Shift+Tab} backward. \texttt{?} toggles help (\texttt{Esc} also closes it).
\texttt{q} or \texttt{Ctrl+C} quits.
\begin{center}\small
\begin{tabularx}{\textwidth}{@{}L{3cm}L{4.2cm}Y@{}}
\toprule
\textbf{Tab} & \textbf{Keys} & \textbf{Meaning}\\
\midrule
Processes & \texttt{j}/\texttt{k} or arrows & Move selection through the process table (PID, COMM, ZONE, SCORE, UID, CONTAINMENT).\\
 & \texttt{/} & Start a filter; type to narrow, \texttt{Enter}/\texttt{Esc} to stop typing, \texttt{Esc} again clears it.\\
 & \texttt{c} & Contain the selected process (see below).\\
Alerts & \texttt{j}/\texttt{k}, \texttt{Enter} & Scroll alerts; \texttt{Enter} expands or collapses one.\\
Agent Activity & \texttt{j}/\texttt{k} & Scroll the swarm's activity feed.\\
Query Console & type, \texttt{Enter} & Run a command; \texttt{Up}/\texttt{Down} history; \texttt{Ctrl+U} clears the line.\\
\bottomrule
\end{tabularx}
\end{center}

\subsection{Containing a process}
\resetsteps
\step{Select it and press \texttt{c}.} A modal asks for a level.
\step{Pick the level.}
\begin{center}\small
\begin{tabular}{@{}ll@{}}
\toprule \textbf{Key} & \textbf{Level}\\ \midrule
\texttt{0} & None (release)\\ \texttt{1} & Cgroup\\ \texttt{2} & Seccomp\\ \texttt{3} & Namespace\\ \texttt{4} & Terminate\\ \bottomrule
\end{tabular}
\end{center}
\step{Confirm.} \texttt{y} or \texttt{Enter} submits with the reason \texttt{operator action via kb-tui};
\texttt{n} returns to level selection; \texttt{Esc} cancels.

\subsection{The query console}
\begin{lstlisting}[style=console]
> help
> pid 4242                      # full state for one process
> zone borderlands              # list processes in safe|suspicious|borderlands
> stats                         # global telemetry snapshot
> contain 4242 seccomp suspicious beaconing
> clear
\end{lstlisting}
Levels for \texttt{contain}: \texttt{none, cgroup, seccomp, namespace, terminate}.

\subsection{When \cmd{kbd} is unreachable}
If \file{/run/kb/kba.sock} cannot be dialed the TUI does \emph{not} fail. It shows an
\texttt{[OFFLINE]} banner (``\texttt{kbd unreachable}'') and a \texttt{[DEMO DATA]} table so you can see the
interface working. Never act on demo rows: the PIDs are fictional. Read the banner before the table.
